\subsection{幺半群的全代数}

幺半群 S 在 A 上的代数（作为 A-模）是积 \(A^{s}\) 中由具有有限支集的族 \((\alpha_{s})_{s\in\mathbb{S}}\) 组成的子模；这个代数中的乘法由关系式 \((\alpha_{s})(\beta_{s})=(\gamma_{s})\) 定义，其中，对所有 \(s\in S\) ，

\[
\gamma_ {s} = \sum_ {t u = s} \alpha_ {t} \beta_ {u}\tag{38}
\]

（参见第 6 号，公式 (35)）。式 (38) 右端的和是有意义的，因为 \((\alpha_{s})\) 和 \((\beta_{s})\) 都是有限支集的族，从而双重族 \((\alpha_{t}\beta_{u})_{(t,u)\in S\times S}\) 也是有限支集的。但式 (38) 的右端对任意元素 \((\alpha_{s})\) , \(\beta_{s})\) （属于 \(A^{S}\) ）也是有意义的，只要幺半群 \(S\) 满足以下条件：

\begin{enumerate}
\def\labelenumi{(\Alph{enumi})}
\setcounter{enumi}{3}
\tightlist
\item
  对所有 \(s \in \mathbf{S}\) ，只存在有限多个有序对 \((t, u)\) —— \(\mathbf{S} \times \mathbf{S}\) 中的——，使得 \(tu = s\) 。
\end{enumerate}

现在假设 S 满足条件 (D)；此时可以在乘积 A-模 \(A^{s}\) 上由公式 (38) 定义一个乘法。立即可见，这样定义在 \(A^{s}\) 上的乘法是 A-双线性的；它也是结合的，因为对于 \(\alpha\) , \(\beta\) , \(\gamma\) （属于 \(A^{s}\) ），

\[
\sum_ {u v w = t} \alpha_ {u} \beta_ {v} \gamma_ {w} = \sum_ {r w = t} \left(\left(\sum_ {u v = r} \alpha_ {u} \beta_ {v}\right) \gamma_ {w}\right) = \sum_ {u s = t} \left(\alpha_ {u} \left(\sum_ {v w = s} \beta_ {v} \gamma_ {w}\right)\right).
\]

该乘法与 \(A^{s}\) 上的 A-模结构一起，在 \(A^{s}\) 上定义了一个 A 上的含单位元结合代数结构；我们称集合 \(A^{s}\) 连同此结构为幺半群 S 在 A 上的全代数。

立即可见，代数 \(\mathbf{A}^{(\mathbf{S})}\) ——幺半群 \(\mathbf{S}\) 在 \(\mathbf{A}\) 上的——（在需要避免混淆时也称 \(\mathbf{S}\) 的限制代数）是幺半群 \(\mathbf{S}\) 在 \(\mathbf{A}\) 上的全代数的一个子代数（且当 \(\mathbf{S}\) 有限时与后者相同）。作为语言上的滥用，全代数的每个元素 \((\xi_s)_{s \in \mathbf{S}}\) （ \(\mathbf{S}\) 在 \(\mathbf{A}\) 上的）也沿用同样的记号 \(\sum_{s \in \mathbf{S}} \xi_s e_s\) （或甚至 \(\sum_{s \in \mathbf{S}} \xi_s s\) ）—— \(\mathbf{S}\) 的限制代数的元素所用的记号——；当然，这个记号中出现的求和符号不对应任何代数运算，因为它一般是对无穷多个 \(\neq 0\) 的项取的。利用这一记号， \(\mathbf{S}\) 的全代数中的乘法也由第 6 号的公式 (35) 给出。

若 S 是交换的，则它的全代数 \(A^{S}\) 也是交换的。若 T 是 S 的子幺半群，则幺半群的全代数 \(A^{T}\) 典范地等同于 S 的全代数的一个子代数。若 \(\rho: A \to B\) 是环同态，则典范扩张 \(\rho^{S}: A^{S} \to B^{S}\) 是从 S 在 A 上的全代数到 S 在 B 上的全代数中的 A-同态，它扩张了典范同态 \(A^{(S)} \to B^{(S)}\) 。
% semantic-fragments: legacy-input-seam
\relax{}
